\documentclass[12pt]{article}
\usepackage[utf8]{inputenc}
\usepackage[T1]{fontenc}
\usepackage{hyperref}
\title{Sukurum Language Resources for K-12 Curriculum Development}
\author{Compiled by Claude Sonnet 5.5 (claude-sonnet-5-5) for Andr\'as Kornai}
\date{}
\begin{document}
\maketitle

\section{Introduction}
Sukurum (ISO 639-3 zsu; Glottocode suku1264; also called Wapu) is a language of Papua New Guinea. The survey description lists it as Austronesian with about 1,200 speakers. Sources found in search disagree on the details. Global Recordings Network gives a population of 1,900. Joshua Project gives 2,400 speakers. Joshua Project's people-group page and its language page also classify the language differently. No school materials were found, and the language is not known to be taught in schools.

\section{Primary Educational Resources}
No primers, textbooks or school programmes for Sukurum turned up in the searches.

\section{Linguistic and Reference Materials}
\begin{itemize}
\item PARADISEC catalogue: Sukurum recordings from 1985 collected by Susanne Holzknecht, with items tagged zsu. \url{https://catalog.paradisec.org.au/repository/SH2/ANUA520_3310}
\item ANU Archives: collection of Sukurum cassette tapes (1--11). \url{https://archivescollection.anu.edu.au/index.php/papua-new-guinea-language-recordings-sukurum-casette-tapes-1-11/informationobject/reports}
\end{itemize}
No dictionary or grammar was found in the search results.

\section{Religious and Community Resources}
\begin{itemize}
\item Joshua Project language page for Sukurum. It lists Bible translation as needed, with no Bible audio or Jesus Film available. \url{https://joshuaproject.net/languages/zsu}
\item Joshua Project people-group page for the Sukurum. \url{https://joshuaproject.net/people_groups/18904}
\end{itemize}

\section{Digital and Audio Resources}
\begin{itemize}
\item Global Recordings Network page for zsu. The language is listed, but no audio recordings were available at the time of the search. \url{https://globalrecordings.net/en/language/zsu}
\end{itemize}

\section{Curriculum Development Context}
Very little published material exists for Sukurum. A curriculum would have to be built from scratch, starting with the archival recordings held at PARADISEC and ANU and with community-led documentation. The orthography and the language's classification should be confirmed with specialists and the community before any materials are produced.

\end{document}